\section{Magisterial Chronology}
\label{app:magisterial-chronology}

This chronological register records the conciliar definitions, creeds, papal constitutions, and curial directives that teach the faith concerning the angels or regulate angelic devotion. Each entry states the act's date, promulgating authority, kind, doctrinal or disciplinary substance, and grade of authority attached to its core proposition.

{\footnotesize
\setlength{\tabcolsep}{2.5pt}
\renewcommand{\arraystretch}{1.15}
\begin{longtable}{>{\bfseries\raggedright\arraybackslash}p{0.10\linewidth}>{\raggedright\arraybackslash}p{0.21\linewidth}>{\raggedright\arraybackslash}p{0.16\linewidth}>{\raggedright\arraybackslash}p{0.34\linewidth}>{\raggedright\arraybackslash}p{0.13\linewidth}}
\toprule
\textbf{Date} & \textbf{Act and Authority} & \textbf{Kind} & \textbf{Teaching or Regulation} & \textbf{Proposition Grade} \\
\midrule
\endfirsthead
\toprule
\textbf{Date} & \textbf{Act and Authority} & \textbf{Kind} & \textbf{Teaching or Regulation} & \textbf{Proposition Grade} \\
\midrule
\endhead
\bottomrule
\endfoot
325 / 381 & First Council of Nicaea \& First Council of Constantinople & Ecumenical creed & Confesses God the Father Almighty as Maker of all things visible and invisible (\emph{visibilium omnium et invisibilium}), encompassing the angelic world. & \textbf{Defined} \\
\addlinespace[2pt]
c.~343--381 & Synod of Laodicea, Canon~35 & Regional synod canon & Forbids Christians from forsaking the Church to name and invoke angels in secret idolatrous assemblies (Hefele, vol.~2). & Disciplinary / \textbf{Church teaching} \\
\addlinespace[2pt]
543 / 553 & Synod of Constantinople (anti-Origenist anathemas) & Conciliar / imperial-synodal anathemas & Condemns the claim that the punishment of demons and wicked men is temporary and will end in universal restoration (\gk{apokatastasis}; anath.~9, Hefele, vol.~4). & \textbf{Defined} (in universal reception) \\
\addlinespace[2pt]
745 & Roman Synod under Pope St.~Zacharias & Papal synod & Condemns the Frankish preacher Aldebert for invoking extra-scriptural angelic names in prayer, restricting liturgical invocation to canonical names (Michael, Gabriel, Raphael; Hefele, vol.~2). & Disciplinary / \textbf{Church teaching} \\
\addlinespace[2pt]
1215 & Fourth Lateran Council, \emph{Firmiter credimus} (DS~800--801) & Ecumenical council definition & Defines that God created from nothing, from the beginning of time, both spiritual (angelic) and corporeal creation; that the devil and other demons were created good by nature and became evil by themselves; and that the wicked receive eternal punishment with the devil. & \textbf{Defined} \\
\addlinespace[2pt]
1312 & Council of Vienne, \emph{Ad nostrum qui} (DS~891--899) & Ecumenical council constitution & Condemns the claim that every intellectual nature is naturally blessed in itself and needs no supernatural light of glory to see and enjoy God. & \textbf{Defined} \\
\addlinespace[2pt]
1439 / 1442 & Council of Florence, \emph{Laetentur caeli} \& \emph{Cantate Domino} (DS~1305, 1333, 1336) & Ecumenical council decrees & Affirms immediate beatific vision of the Triune God and reaffirms against Manichean dualism that no nature is evil because God created all things good, including spiritual creatures (\emph{Bullarium Romanum}, vol.~V). & \textbf{Defined} \\
\addlinespace[2pt]
1546--1563 & Council of Trent, Sess.~5, 6, 25 \& \emph{Roman Catechism} (1566) & Ecumenical council \& official catechism & Affirms demonic captivity and deliverance through Christ (Sess.~5, 6); vindicates legitimate veneration and iconography of angels; teaches creation of angels and personal angelic guardianship (\emph{Roman Catechism}, Creed Art.~I, 1st Commandment, Lord's Prayer). & \textbf{Defined} (redemption) / \textbf{Church teaching} (guardianship) \\
\addlinespace[2pt]
1870 & First Vatican Council, \emph{Dei Filius}, Ch.~1 \& Can.~I.1--5 & Ecumenical council dogmatic constitution & Reaffirms verbatim Lateran~IV's definition that God freely created out of nothing, from the beginning of time, both spiritual (angelic) and corporeal creatures for His glory. & \textbf{Defined} \\
\addlinespace[2pt]
1950 & Pope Pius~XII, \emph{Humani generis}, \S~26 & Papal encyclical & Reproves questioning whether angels are personal beings, upholding their personal, spiritual reality. & \textbf{Church teaching} \\
\addlinespace[2pt]
1954 & Pope Pius~XII, \emph{Ad caeli Reginam} & Papal encyclical & Teaches the royal dignity of the Blessed Virgin Mary as Queen of Heaven and of the angels. & \textbf{Church teaching} \\
\addlinespace[2pt]
1964 & Second Vatican Council, \emph{Lumen gentium}, \S\S~53, 56, 59, 66 & Ecumenical council dogmatic constitution & Teaches that Mary, though redeemed by her Son, surpasses all creatures in heaven and on earth (including the angels) and receives veneration (\emph{dulia}/\emph{hyperdulia}) strictly distinct from adoration (\emph{latria}). & \textbf{Church teaching} \\
\addlinespace[2pt]
1975 & Congregation for the Doctrine of the Faith, \emph{Christian Faith and Demonology} & Curial doctrinal study & Reviews biblical, patristic, and conciliar doctrine on Satan and the demons, reaffirming personal demonic agency created good and fallen freely. & \textbf{Church teaching} \\
\addlinespace[2pt]
1992 & Congregation for the Doctrine of the Faith, Decree on \emph{Opus Angelorum} & Curial disciplinary decree & Regulates devotion to the angels, forbidding the use of extra-scriptural angelic names derived from private revelations in prayers and liturgical practices. & Disciplinary / \textbf{Church teaching} \\
\addlinespace[2pt]
1992 / 1997 & \emph{Catechism of the Catholic Church}, \S\S~325--336, 391--395 & Universal catechism & Synthesizes Catholic teaching on angels as purely spiritual, personal, immortal creatures, Christ's servants, personal guardian angels (\S~336), and the irrevocable fall and bounded power of Satan (\S\S~391--395). & \textbf{Church teaching} (incorporating \textbf{Defined} articles) \\
\addlinespace[2pt]
2001 & Congregation for Divine Worship, \emph{Directory on Popular Piety and the Liturgy}, \S\S~212--217 & Curial pastoral directory & Guides legitimate devotion to the holy angels while discouraging assigning proper names to angels other than Michael, Gabriel, and Raphael. & Disciplinary / \textbf{Church teaching} \\
\end{longtable}
}
